% ID: FIS-TER-GAS-005
% Titolo: Gas monoatomico in una bombola rigida
% Disciplina: Fisica
% Area: Termologia
% Argomento: Gas perfetti
% Sottoargomento: Trasformazione isocora, primo principio ed energia interna
% Classe: Quarta liceo scientifico
% Difficolta: 3
% Tipo: problema_numerico
% Risultato: Q = 252 J; Delta U = 252 J
% Tempo_stimato: 10 min
% Competenze: applicazione dell'equazione di stato, primo principio della termodinamica, energia interna di un gas monoatomico
% Prerequisiti: gas perfetti, trasformazioni isocore, primo principio della termodinamica, energia interna
% Tag: termologia, gas_perfetti, trasformazione_isocora, primo_principio, energia_interna, gas_monoatomico
% Fonte: esercizio_proposto_utente
% Autore: Riccardo Carli
% Licenza: CC-BY-SA-4.0
% Simulazione: ideal_gas_process

\begin{esercizio}
In una bombola e contenuto un volume di
\[
V=6{,}00\cdot 10^{-2}\,\mathrm{m^3}
\]
di un gas perfetto monoatomico. In seguito a una trasformazione termodinamica,
la pressione del gas aumenta di \(2800\,\mathrm{Pa}\).

Determina:
\begin{enumerate}
    \item quanto calore ha assorbito il gas;
    \item qual e la sua variazione di energia interna.
\end{enumerate}
\end{esercizio}

\begin{soluzione}
Poiche il gas e contenuto in una bombola rigida, il volume resta costante. La
trasformazione e quindi isocora.

Per un gas perfetto vale l'equazione di stato
\[
pV=nRT.
\]
Scrivendola per lo stato iniziale e per quello finale e sottraendo membro a
membro, dato che \(V\) e costante, si ottiene
\[
(p_f-p_i)V=nR(T_f-T_i),
\]
cioe
\[
V\,\Delta p=nR\,\Delta T.
\]

Con i dati del problema:
\[
nR\,\Delta T=V\,\Delta p
=(6{,}00\cdot10^{-2})(2800)=168\,\mathrm{J}.
\]

Per un gas perfetto monoatomico l'energia interna e
\[
U=\frac{3}{2}nRT,
\]
percio la sua variazione e
\[
\Delta U=\frac{3}{2}nR\,\Delta T.
\]
Usando la relazione ricavata sopra:
\[
\Delta U=\frac{3}{2}V\,\Delta p
=\frac{3}{2}(168\,\mathrm{J})
=252\,\mathrm{J}.
\]
Quindi
\[
\boxed{\Delta U=252\,\mathrm{J}}.
\]

Applichiamo ora il primo principio della termodinamica nella convenzione
\[
\Delta U=Q-L,
\]
dove \(L\) e il lavoro compiuto dal gas. In una trasformazione isocora
\(\Delta V=0\), quindi
\[
L=0.
\]
Ne segue che
\[
Q=\Delta U,
\]
e dunque
\[
\boxed{Q=252\,\mathrm{J}}.
\]
Il segno positivo indica che il gas ha assorbito calore.

\medskip
\textbf{Approfondimento sul numero di moli.}
Con i soli dati forniti dal problema non e possibile determinare numericamente
il numero di moli \(n\). Infatti dalla trasformazione isocora si ricava
\[
n=\frac{V\,\Delta p}{R\,\Delta T},
\]
ma il valore di \(\Delta T\) non e noto. Equivalentemente, usando l'equazione
di stato \(n=pV/(RT)\), sarebbero necessari la pressione assoluta e la
temperatura di almeno uno dei due stati. Il problema e risolvibile senza
conoscere \(n\) proprio perche il prodotto \(nR\,\Delta T\) puo essere
sostituito con \(V\,\Delta p\).
\end{soluzione}
